\chapter{The Transformer architecture}
\label{ch:transformers}

\begin{quote}
\emph{``Attention is all you need.'' --- Vaswani et al., 2017}
\end{quote}

% ============================================================
\section{From RNNs to Transformers}

Recurrent neural networks (RNNs, LSTMs) process sequences token by token, left to right. This creates two problems:
\begin{enumerate}
  \item \textbf{Sequential bottleneck.} Token $t$ must wait for tokens $1, \ldots, t-1$ to be processed. No parallelism.
  \item \textbf{Vanishing context.} Information from early tokens fades as the sequence grows.
\end{enumerate}

The Transformer solves both problems with \emph{attention}: every token can attend to every other token directly, in parallel.

% ============================================================
\section{Scaled dot-product attention}

Given a sequence of $n$ tokens, each represented by a vector, we compute three matrices:
\begin{itemize}
  \item $\mathbf{Q}$ (queries), $\mathbf{K}$ (keys), $\mathbf{V}$ (values), each of shape $(n \times d_k)$.
\end{itemize}

\[
\text{Attention}(\mathbf{Q}, \mathbf{K}, \mathbf{V}) = \text{softmax}\!\left(\frac{\mathbf{Q}\mathbf{K}^\top}{\sqrt{d_k}}\right)\mathbf{V}
\]

The $\sqrt{d_k}$ scaling prevents the dot products from growing too large, which would push softmax into regions with vanishing gradients.

\begin{minted}{python}
import torch
import torch.nn.functional as F

def scaled_dot_product_attention(Q, K, V, mask=None):
    """Compute scaled dot-product attention."""
    d_k = Q.size(-1)
    scores = torch.matmul(Q, K.transpose(-2, -1)) / (d_k ** 0.5)
    if mask is not None:
        scores = scores.masked_fill(mask == 0, float("-inf"))
    weights = F.softmax(scores, dim=-1)
    return torch.matmul(weights, V), weights

# Example: 1 batch, 4 tokens, embedding dim 8
Q = torch.randn(1, 4, 8)
K = torch.randn(1, 4, 8)
V = torch.randn(1, 4, 8)
output, attn_weights = scaled_dot_product_attention(Q, K, V)
print(f"Output shape: {output.shape}")       # (1, 4, 8)
print(f"Attention weights shape: {attn_weights.shape}")  # (1, 4, 4)
\end{minted}

\begin{aitip}
The attention weight matrix is $(n \times n)$. Entry $(i, j)$ tells you how much token $i$ ``looks at'' token $j$. Visualizing this matrix reveals what the model focuses on.
\end{aitip}

% ============================================================
\section{Multi-head attention}

Instead of one attention function, Transformers use $h$ parallel ``heads,'' each with its own learned projections $W_Q^{(i)}, W_K^{(i)}, W_V^{(i)}$:

\[
\text{MultiHead}(\mathbf{Q}, \mathbf{K}, \mathbf{V}) = \text{Concat}(\text{head}_1, \ldots, \text{head}_h)\,\mathbf{W}^O
\]

Each head can learn a different type of relationship: syntactic, semantic, positional, etc.

\begin{minted}{python}
import torch.nn as nn

class MultiHeadAttention(nn.Module):
    def __init__(self, d_model=512, num_heads=8):
        super().__init__()
        assert d_model % num_heads == 0
        self.d_k = d_model // num_heads
        self.num_heads = num_heads
        self.W_q = nn.Linear(d_model, d_model)
        self.W_k = nn.Linear(d_model, d_model)
        self.W_v = nn.Linear(d_model, d_model)
        self.W_o = nn.Linear(d_model, d_model)

    def forward(self, Q, K, V, mask=None):
        batch = Q.size(0)
        # Project and reshape: (batch, seq, d_model) -> (batch, heads, seq, d_k)
        Q = self.W_q(Q).view(batch, -1, self.num_heads, self.d_k).transpose(1, 2)
        K = self.W_k(K).view(batch, -1, self.num_heads, self.d_k).transpose(1, 2)
        V = self.W_v(V).view(batch, -1, self.num_heads, self.d_k).transpose(1, 2)
        # Attention per head
        out, weights = scaled_dot_product_attention(Q, K, V, mask)
        # Concatenate heads
        out = out.transpose(1, 2).contiguous().view(batch, -1, self.num_heads * self.d_k)
        return self.W_o(out)

mha = MultiHeadAttention(d_model=64, num_heads=4)
x = torch.randn(2, 10, 64)  # batch=2, seq=10, dim=64
print(f"Output: {mha(x, x, x).shape}")  # (2, 10, 64)
\end{minted}

% ============================================================
\section{Positional encoding}

Attention is permutation-invariant --- it does not know token order. Positional encodings inject position information:

\[
PE_{(pos, 2i)} = \sin\!\left(\frac{pos}{10000^{2i/d}}\right), \quad
PE_{(pos, 2i+1)} = \cos\!\left(\frac{pos}{10000^{2i/d}}\right)
\]

\begin{minted}{python}
import numpy as np
import matplotlib.pyplot as plt

def positional_encoding(max_len, d_model):
    pe = np.zeros((max_len, d_model))
    position = np.arange(max_len)[:, np.newaxis]
    div_term = np.exp(np.arange(0, d_model, 2) * -(np.log(10000.0) / d_model))
    pe[:, 0::2] = np.sin(position * div_term)
    pe[:, 1::2] = np.cos(position * div_term)
    return pe

pe = positional_encoding(100, 64)
plt.figure(figsize=(10, 4))
plt.imshow(pe, aspect="auto", cmap="RdBu")
plt.xlabel("Embedding dimension")
plt.ylabel("Position")
plt.title("Sinusoidal positional encoding")
plt.colorbar()
plt.tight_layout()
plt.savefig("positional_encoding.png", dpi=150)
plt.show()
\end{minted}

\begin{aitip}
Modern models (GPT, Llama) use \emph{Rotary Position Embeddings} (RoPE) instead of sinusoidal encodings. RoPE encodes relative position by rotating query and key vectors, which scales better to long contexts.
\end{aitip}

% ============================================================
\section{Encoder-decoder architecture}

The original Transformer has two stacks:
\begin{itemize}
  \item \textbf{Encoder:} processes the input sequence with bidirectional self-attention. Each token can see all other tokens.
  \item \textbf{Decoder:} generates the output sequence with causal (masked) self-attention. Each token can only see previous tokens. It also attends to the encoder output via cross-attention.
\end{itemize}

\begin{center}
\begin{tikzpicture}[
  block/.style={rectangle, draw, fill=aiblue!15, minimum width=3cm, minimum height=0.8cm, rounded corners},
  arrow/.style={-{Stealth}, thick}
]
  \node[block] (emb_enc) at (0, 0) {Input Embedding + PE};
  \node[block] (sa_enc) at (0, 1.2) {Multi-Head Self-Attention};
  \node[block] (ff_enc) at (0, 2.4) {Feed-Forward Network};
  \node[above] at (0, 3.1) {\textbf{Encoder} ($\times N$)};

  \node[block] (emb_dec) at (5.5, 0) {Output Embedding + PE};
  \node[block] (sa_dec) at (5.5, 1.2) {Masked Self-Attention};
  \node[block] (ca_dec) at (5.5, 2.4) {Cross-Attention};
  \node[block] (ff_dec) at (5.5, 3.6) {Feed-Forward Network};
  \node[block] (linear) at (5.5, 4.8) {Linear + Softmax};
  \node[above] at (5.5, 5.5) {\textbf{Decoder} ($\times N$)};

  \draw[arrow] (emb_enc) -- (sa_enc);
  \draw[arrow] (sa_enc) -- (ff_enc);
  \draw[arrow] (emb_dec) -- (sa_dec);
  \draw[arrow] (sa_dec) -- (ca_dec);
  \draw[arrow] (ca_dec) -- (ff_dec);
  \draw[arrow] (ff_dec) -- (linear);
  \draw[arrow, dashed] (ff_enc.east) -- ++(1,0) |- (ca_dec.west);
\end{tikzpicture}
\end{center}

% ============================================================
\section{Variants in practice}

\begin{longtable}{p{3.5cm}p{3.5cm}p{5cm}}
\toprule
\textbf{Architecture} & \textbf{Models} & \textbf{Use case} \\
\midrule
Encoder-only & BERT, DistilBERT, RoBERTa & Classification, NER, embeddings \\
Decoder-only & GPT-2, GPT-4, Llama, Phi & Text generation, chatbots, code \\
Encoder-decoder & T5, BART, Flan-T5 & Translation, summarization \\
\bottomrule
\end{longtable}

% ============================================================
\section{Visualizing attention with DistilBERT}

\begin{minted}{python}
from transformers import AutoTokenizer, AutoModel
import torch
import matplotlib.pyplot as plt
import seaborn as sns

model_name = "distilbert-base-uncased"
tokenizer = AutoTokenizer.from_pretrained(model_name)
model = AutoModel.from_pretrained(model_name, output_attentions=True)

text = "The cat sat on the mat because it was tired"
inputs = tokenizer(text, return_tensors="pt")
with torch.no_grad():
    outputs = model(**inputs)

# outputs.attentions is a tuple: one tensor per layer
# Each tensor shape: (batch, heads, seq_len, seq_len)
attn = outputs.attentions[-1][0]  # last layer, first batch
tokens_list = tokenizer.convert_ids_to_tokens(inputs["input_ids"][0])

fig, axes = plt.subplots(1, 4, figsize=(20, 5))
for i, ax in enumerate(axes):
    sns.heatmap(attn[i].numpy(), xticklabels=tokens_list,
                yticklabels=tokens_list, ax=ax, cmap="Blues")
    ax.set_title(f"Head {i}")
plt.tight_layout()
plt.savefig("attention_heads.png", dpi=150)
plt.show()
\end{minted}

\begin{exercise}
\begin{enumerate}
  \item Implement the \texttt{scaled\_dot\_product\_attention} function from scratch (no PyTorch attention API). Verify it matches the output of \texttt{torch.nn.functional.scaled\_dot\_product\_attention}.
  \item Visualize attention heads for the sentence ``The bank by the river was steep.'' Which heads attend to which words? Do any heads capture the word ``bank'' relating to ``river''?
  \item Modify the positional encoding function to use learned embeddings instead of sinusoidal. Train a small model and compare.
  \item Compute the number of parameters in a Transformer encoder layer with $d_{\text{model}}=512$ and $h=8$. Show your calculation.
\end{enumerate}
\end{exercise}

\begin{ethicsbox}
Attention weights are sometimes interpreted as ``explanations'' of model behavior. This is misleading. Attention shows where the model \emph{looks}, not necessarily what it \emph{uses} for the final prediction. Use attention visualization as a diagnostic tool, not as a faithful explanation.
\end{ethicsbox}

% ============================================================
\section{Chapter summary}

\begin{itemize}
  \item The Transformer replaces recurrence with self-attention, enabling full parallelism.
  \item Scaled dot-product attention computes a weighted sum of values, with weights from query-key similarity.
  \item Multi-head attention allows the model to attend to different relationship types simultaneously.
  \item Positional encodings (sinusoidal or learned) inject sequence order into the attention mechanism.
  \item Encoder-only (BERT), decoder-only (GPT), and encoder-decoder (T5) are the three main Transformer variants.
  \item Attention visualization is useful for debugging but should not be over-interpreted as model explanation.
\end{itemize}
